\documentclass[11pt, a4paper]{article}
\usepackage[a4paper, margin=2.5cm]{geometry}
\usepackage[utf8]{inputenc}
\usepackage[spanish]{babel}
\usepackage{booktabs}
\usepackage{tabularx}
\usepackage{xcolor}
\usepackage{colortbl} % <-- Required for \rowcolor support

\begin{document}

\noindent
\textbf{Nivel Diagnóstico: Estudio de Mapeo Sistemático (SMS - Petersen et al.)}

\vspace{0.3cm}

\noindent
\begin{tabularx}{\textwidth}{|X|X|}
\hline
\rowcolor{gray!15}
\textbf{Preguntas de investigación del relevamiento (RQs)} & \textbf{Protocolo de extracción y clasificación} \\ \hline
\textbf{RQ1:} ¿Qué sistemas o plataformas tecnológicas se han propuesto en la literatura para gestionar pasantías y prácticas profesionales entre universidades, empresas y estudiantes?
\vspace{0.2cm}

\textbf{RQ2:} ¿Qué enfoques de diseño, métodos o arquitecturas utilizan con más frecuencia esas soluciones?
\vspace{0.2cm}

\textbf{RQ3:} qué limitaciones o vacíos de investigación (\textit{research gaps}) siguen sin resolverse en las soluciones existentes? & 
\textbf{Estrategia de búsqueda:} se buscará en bases de datos académicas y en repositorios de tesis universitarias, en español y en inglés, con cadenas de búsqueda que combinen.
\vspace{0.2cm}

\textbf{Selección (inclusión y exclusión):} se revisarán títulos, resúmenes y palabras clave, y el texto completo cuando haya dudas. Se incluirán los estudios que presenten una solución tecnológica para gestionar pasantías o prácticas, publicados desde 2015, y se excluirán los duplicados y los documentos sin texto completo.
\vspace{0.2cm}

\textbf{Clasificación (\textit{keywording}):} cada estudio se agrupará según el tipo de solución (ofertas, postulación, seguimiento de horas, comunicación universidad-empresa), el tipo de aporte (método, herramienta o modelo) y el tipo de investigación (propuesta de solución o evaluación).
\vspace{0.2cm}

\textbf{Extracción y mapeo:} los datos de cada estudio (año, país, funciones, arquitectura y limitaciones) se registrarán en una matriz. Con ella se elaborará un gráfico de burbujas (\textit{bubble plot}) que muestre cuántos estudios hay por categoría e identifique los temas poco explorados. \\ \hline
\end{tabularx}

\end{document}